\documentclass{article}
\usepackage[utf8]{inputenc}

\usepackage{mathpartir}
\usepackage{syntax}
\usepackage{amssymb}

\newcommand{\irule}[3]{ \inferrule*[right=(#1)]{#2}{#3} }

% Fix '::' -> ': :' when using babel
\newenvironment{grammarf}{\shorthandoff{:}\begin{grammar}}{\end{grammar}\shorthandon{:}}

\newcommand{\gramisa}[0]{::=~}
\newcommand{\gramor}[0]{|\hspace{3mm}}

\newcommand{\gramlr}[2]{%
  \begin{minipage}{0.2\textwidth}\raggedleft #1\end{minipage}
  \begin{minipage}{0.8\textwidth}#2\end{minipage}
}

% Types
\newcommand{\tInt}{\textsc{Int}\xspace}
\newcommand{\tFloat}{\textsc{Float}\xspace}
\newcommand{\tVoid}{\textsc{Void}\xspace}

\newcommand{\tFV}[1]{\textrm{FreeVars}(#1)}

\def\opxor{\mathbin{\char`\^}}

% Constant
\newcommand{\cNil}{\textsc{Nil}}

% Program & language names
\newcommand{\ctonewspeak}{\textsc{c2newspeak}}
\newcommand{\cpp}{\textsc{cpp}}
\newcommand{\newspeak}{\textsc{Newspeak}}

% Newspeak
\newcommand{\npkDecl}[2]{\textsc{Decl}~(#1) \{ #2 \}}
\newcommand{\npkIf}[3]{\textsc{If}(#1) \{ #2 \} \textsc{Else} \{ #3 \}}
\newcommand{\npkDoWith}[2]{\textsc{Do}\{ #1 \}\textsc{With}~#2:}
\newcommand{\npkWhile}[1]{\textsc{While}(1)\{#1\}}
\newcommand{\npkGoto}[1]{\textsc{Goto}~#1}
\newcommand{\npkReturn}[1]{\textsc{Return}~#1}

% Continuations
\def\kPass{\textsc{Kpass}}
\def\kSeq#1#2{\textsc{Kseq}(#1, #2)}
\def\kBlock#1#2{\textsc{Kblock}(#1, #2)}
\def\kGoto#1{\textsc{Kgoto}(#1)}
\def\kCall#1#2#3#4{\textsc{Kcall}(#1, #2, #3, #4)}
\def\kRet#1{\textsc{Kret}(#1)}

\newcommand{\insertcode}[1]{\input{gen/#1.pyg.tex}}

% \chapter* + add to toc
%\newcommand{\chapterstar}[1]{\chapter*{#1} \addcontentsline{toc}{chapter}{#1}}

% Register names
\def\defreg#1{\expandafter\def\csname#1\endcsname{\textsc{#1}\xspace}}

\defreg{eax}
\defreg{ebx}
\defreg{ecx}
\defreg{edx}
\defreg{esi}
\defreg{edi}
\defreg{ebp}
\defreg{esp}
\defreg{eip}
\defreg{eflags}

\def\crtrois{\textsc{CR3}}

\newcommand{\ring}[1]{\emph{ring}~#1}

\newcommand{\asminstr}[1]{\texttt{#1}}

\newcommand{\wip}[0]{%
\begin{center}
    \includegraphics[width=5cm]{gen/travaux.pdf}
\end{center}
}

\newcommand{\listcons}[0]{%
\texttt{(::)}%
}

\newcommand{\hexa}[1]{%
\texttt{0x#1}%
}

% Debug tikz bounding box
\def\drawbb{\draw[thick] (current bounding box.south west) rectangle (current bounding box.north east);}

% comes from the memoir manual {{{

\newcommand*{\addtotextheight}[1]{%
\twocolumn
\addtolength{\textheight}{#1}%
\onecolumn}

%% for main body, bottom of text at 1in, footer below
%% top of header at 1in, first text line double spaced
%% below base of header
\newlength{\linespace}
\setlength{\linespace}{\baselineskip} % current equivalent of \onelineskip
\setlength{\headheight}{\onelineskip}
\setlength{\headsep}{\linespace}
\addtolength{\headsep}{-\topskip}
\setlength{\uppermargin}{1in}
\addtolength{\uppermargin}{\headheight}
\addtolength{\uppermargin}{\headsep}

%% and for the bottom
\setlength{\lowermargin}{1in}
\setlength{\textheight}{\paperheight}

\addtolength{\textheight}{-\uppermargin}
\addtolength{\textheight}{-\lowermargin}
%% footnote settings
\setlength{\footskip}{\onelineskip}
\setlength{\footnotesep}{\onelineskip}


%% the fiddle lengths (..ta.. for title/approval page, others for prelims)
\newlength{\toptafiddle} \setlength{\toptafiddle}{2\linespace}
\newlength{\bottafiddle} \setlength{\bottafiddle}{0pt}
\newlength{\topfiddle}
\setlength{\topfiddle}{\toptafiddle}
\newlength{\botfiddle}
\setlength{\botfiddle}{\onelineskip}

%%% general macro for Abstract, etc., headings
\newcommand*{\pretoctitle}[1]{{\clearpage\centering \vspace*{-\topfiddle}#1\par}}

%%% Start the ACKNOWLEDGEMENTS
\newcommand{\asuacknowledgements}{\pretoctitle{ACKNOWLEDGEMENTS}}

%%% Start the ABSTRACT
\newcommand{\asuabstract}{ \addtotextheight{-\botfiddle} \pretoctitle{ABSTRACT}}

\addtodef{\mainmatter}{\addtotextheight{\botfiddle}{}}

%% make it easy to center any dedication
\newcommand{\asudedication}[1]{%
{\clearpage\mbox{}\vfill\centering #1 \par\vfill\clearpage}}

% }}} /memoir

\DeclareUnicodeCharacter{0393}{\Gamma}        % Γ
\DeclareUnicodeCharacter{03B5}{\epsilon}      % ε
\DeclareUnicodeCharacter{03C3}{\sigma}        % σ
\DeclareUnicodeCharacter{2190}{\leftarrow}    % ←
\DeclareUnicodeCharacter{2191}{\uparrow}      % ↑
\DeclareUnicodeCharacter{2192}{\rightarrow}   % →
\DeclareUnicodeCharacter{21D2}{\Rightarrow}   % ⇒
\DeclareUnicodeCharacter{2205}{\emptyset}     % ∅
\DeclareUnicodeCharacter{2208}{\in}           % ∈
\DeclareUnicodeCharacter{2260}{\neq}          % ≠
\DeclareUnicodeCharacter{2295}{\oplus}        % ⊕
\DeclareUnicodeCharacter{22A2}{\vdash}        % ⊢


\begin{document}

\section{Syntaxe}

\begin{minipage}{0.6\textwidth}
\begin{grammar}
<programme> ::= ( fonctions , globales , <bloc> )

    <expr> ::= <lval>
          \alt unop <expr>
          \alt <expr> binop <expr>
          \alt cst
          \alt \& <lval>
          \alt \& fonction

    <lval> ::= var
          \alt <lval> . champ
          \alt <lval> [ <expr> ]
          \alt * <expr>
\end{grammar}
\end{minipage}
\begin{minipage}{0.4\textwidth}
\begin{grammar}
     <bloc> ::= <instr> ; <bloc>
           \alt $ε$

    <instr> ::= <lval> $←$ <expr>
           \alt <lval> $←$ funexp (args)
           \alt  $∅ ←$ funexp (args)
           \alt $↑$ nom \{ <bloc> \}
           \alt if (<expr>) \{ <bloc> \}
           \alt \{ <bloc> \} label:
           \alt goto label
           \alt forever \{ <bloc> \}

\end{grammar}
\end{minipage}

\section{Sémantique}

Sémantique concrète : un état est constitué d'un point de contrôle $p$ et d'un
état mémoire $σ$.

\subsection{Valeurs}

Les valeurs peuvent être :

\begin{itemize}
\item $n$ : un entier
\item $\&a$ : un pointeur vers l'adresse $a$
\item $\&f$ : un pointeur vers la fonction $f$
\end{itemize}

\subsection{Jugements}

Les jugements ont les formes suivantes :

\begin{itemize}

\item $σ ⊢ lv ⇒ a$ :
  la left-value lv correspond à l'adresse mémoire a.

\item $σ ⊢ e ⇒ v$ :
  l'expression e s'évalue en v.

\item $<l, instr, l'>$ :
  on peut passer du point l au point l' en effectuant
  l'instruction instr.

\item $<l, cond, l'>$ :
  on peut passer du point l au point l' si cond est vraie.

\item $(l, σ) → (l', σ')$ :
  permet de définir la fonction de transition principale

\end{itemize}

\subsection{Sémantique des left-values}

\begin{eqnarray*}
\irule{Eval-Lv-Var}{
  (v, a) ∈ σ
}{
  σ ⊢ v ⇒ a
}
&
\irule{Eval-Lv-Deref}{
  σ ⊢ e ⇒ \&a
}{
  σ ⊢ *e ⇒ a
}
\\
\irule{Eval-Lv-Field}{
  σ ⊢ lv ⇒ a
}{
  σ ⊢ lv.f ⇒ a + f
}
&
\irule{Eval-Lv-Array}{
  σ ⊢ lv ⇒ a \\
  σ ⊢ e ⇒ n \\
}{
  σ ⊢ lv[e] ⇒ a + n
}
\end{eqnarray*}

\subsection{Sémantique des expressions}

\begin{eqnarray*}
\irule{Eval-Cst}{
}{
  σ ⊢ c ⇒ c
}
&
\irule{Eval-Lv}{
  σ ⊢ lv ⇒ a \\
  (a, v) ∈ σ
}{
  σ ⊢ lv ⇒ v
}
\\
\irule{Eval-Unop}{
  σ ⊢ e ⇒ v
}{
  σ ⊢ op e ⇒ \widehat{op} v
}
&
\irule{Eval-Binop}{
  σ ⊢ e_1 ⇒ v_1 \\
  σ ⊢ e_2 ⇒ v_2
}{
  σ ⊢ e_1 op e_2 ⇒ v_1 \widehat{op} v_2
}
\\
\irule{Eval-AddrOf}{
  σ ⊢ lv ⇒ a
}{
  σ ⊢ \&lv ⇒ \&a
}
&
\irule{Eval-AddrOfFun}{
}{
  σ ⊢ \&f ⇒ \&f
}
\end{eqnarray*}

\subsection{Sémantique des instructions}

\begin{eqnarray*}
\irule{Instr-Assign}{
  \langle l, lv ← e, l' \rangle \\
  σ ⊢ lv ⇒ a \\
  σ ⊢ e ⇒ v
}{
  (l, σ) → (l', σ ⊕ \{ a → v \})
}
\\
\irule{Instr-Decl}{
  \langle l, ↑x \{ b \}, l' \rangle \\
  σ' = σ ⊕ \{ x → \top \} \\
  \langle l_b, b, l_b' \rangle \\
  (l_b, σ') → (l_b', σ'')
}{
  (l, σ) → (l', σ'' \backslash x)
}
\\
\irule{Instr-Fcall}{
  \langle l, lv ← fe(\vec{e}), l' \rangle \\
  σ ⊢ fe ⇒ f \\
  σ ⊢ \vec{e} ⇒ \vec{v} \\
  σ' = σ ⊕ formals(f, \vec{v}) ⊕ \{ !ret → \top \} \\
  (Entry(f), σ') → (Exit(f), σ'') \\
  σ'' ⊢ !ret ⇒ v_{ret} \\
  σ'' ⊢ lv ⇒ a
}{
  (l, σ) → (l', σ'' \backslash (formals(f) \cup \{!ret\}) ⊕ \{ a → v_{ret}\})
}
\\
\irule{Instr-Fcall-Void}{
  \langle l, ∅ ← fe(\vec{e}), l' \rangle \\
  σ ⊢ fe ⇒ f \\
  σ ⊢ \vec{e} ⇒ \vec{v} \\
  σ' = σ ⊕ formals(f, \vec{v}) \\
  (Entry(f), σ') → (Exit(f), σ'')
}{
  (l, σ) → (l', σ'' \backslash formals(f))
}
\end{eqnarray*}

TODO definir $σ ⊢ fe ⇒ f$

\subsection{Sémantique des conditions}

\[ \irule{Cond}{
  \langle l, e, l' \rangle \\
  σ ⊢ e ⇒ n \\
  n ≠ 0
}{
  (l, σ) → (l', σ)
} \]

\section{Typage}

\subsection{Jugements}

\begin{itemize}
\item $Γ ⊢ b$ :
  sous $Γ$, $b$ est bien typé.
\item $Γ ⊢ e : τ$ :
  sous $Γ$, $e$ est typable en $τ$.
\end{itemize}

\subsection{Flôt de contrôle}

\begin{eqnarray*}
\irule{Pass}{
}{
  Γ ⊢ ε
}
&
\irule{Seq}{
  Γ ⊢ s \\
  Γ ⊢ b
}{
  Γ ⊢ s;b
}
\\
\irule{Forever}{
  Γ ⊢ b
}{
  Γ ⊢ forever \{ b \}
}
&
\irule{If}{
  Γ ⊢ e : Int \\
  Γ ⊢ b
}{
  Γ ⊢ if (e) \{ b \}
} \\
\irule{Goto}{
}{
  Γ ⊢ goto l
}
&
\irule{DoWith}{
  Γ ⊢ b
}{
  Γ ⊢ \{ b \} l:
}
\\
\irule{Decl}{
  Γ ⊕ {x : τ} ⊢ b
}{
  Γ ⊢ ↑x \{b\}
}
&
\irule{Assign}{
  Γ ⊢ lv : τ \\
  Γ ⊢ e  : τ
}{
  Γ ⊢ lv ← e
}
\end{eqnarray*}

\subsection{Expressions}
\begin{eqnarray*}
\irule{Const-Int}{
}{
  Γ ⊢ n : Int
}
&
\irule{Const-Float}{
}{
  Γ ⊢ f : Float
}
\\
\irule{Op-Int}{
  op ∈ \{+,-,×,/,\&,|,\&\&,||,⋘,⋙\} \\
  Γ ⊢ e_1 : Int \\
  Γ ⊢ e_2 : Int
}{
  Γ ⊢ e_1 op e_2 : Int
}
\\
\irule{Op-Float}{
  op ∈ \{+,-,×,/\} \\
  Γ ⊢ e_1 : Float \\
  Γ ⊢ e_2 : Float
}{
  Γ ⊢ e_1 op e_2 : Float
}
\\
\irule{Op-Cmp}{
  op ∈ \{=,≠,≤,≥,<,>\} \\
  Γ ⊢ e_1 : Int \\
  Γ ⊢ e_2 : Int
}{
  Γ ⊢ e_1 op e_2 : Int
}
\\
\irule{Unop-Minus}{
  τ ∈ \{ Int, Float \} \\
  Γ ⊢ e : τ
}{
  Γ ⊢ -e : τ
}
&
\irule{Unop-Not}{
  op ∈ \{ \tilde, ! \}
  Γ ⊢ e : Int
}{
  Γ ⊢ op e : Int
}
\end{eqnarray*}

\end{document}
